\documentclass[12pt, letterpaper]{article}
\usepackage[utf8]{inputenc}
\usepackage{amsmath}
\usepackage{amsfonts}
\usepackage[spanish]{babel}
\usepackage[margin=2.5cm]{geometry}
\usepackage{tikz}
\usetikzlibrary{optics, positioning, arrows.meta, babel}
\usepackage{listings}
\usepackage{xcolor}

\definecolor{codegreen}{rgb}{0,0.6,0}
\definecolor{codegray}{rgb}{0.5,0.5,0.5}
\definecolor{codepurple}{rgb}{0.58,0,0.82}
\definecolor{backcolour}{rgb}{0.95,0.95,0.92}

\lstdefinestyle{mystyle}{
	backgroundcolor=\color{backcolour},   
	commentstyle=\color{codegreen},
	keywordstyle=\color{magenta},
	numberstyle=\tiny\color{codegray},
	stringstyle=\color{codepurple},
	basicstyle=\ttfamily\footnotesize,
	breakatwhitespace=false,         
	breaklines=true,                 
	captionpos=b,                    
	keepspaces=true,                 
	numbers=left,                    
	numbersep=5pt,                  
	showspaces=false,                
	showstringspaces=false,
	showtabs=false,                  
	tabsize=4
}
\lstset{style=mystyle}

\begin{document}
	
	
	$$  I_{dc} = I_{T} = I_O + I_R  \text{ corriente directa / de fondo, porque no lleva información de la fase }$$
	
	$$ I_{ac} = 2\sqrt{I_OI_R} \text{ la amplitud de la parte de la inferferencia } $$
	
	$$ \gamma = \frac{2\sqrt{I_OI_R}}{I_O + I_R}  \text{ es el contraste de la imagen } $$
	
	El $\alpha(t)$ es agregar un valor a la fase, experimentalmente se pone una placa de media o cuarto de onda en el haz objeto o haz de referencia. Tambien se le puede poner un pizoelectrico que dependiendo de la corriente que se aplique va a cambiar la posición.
	
	
	\section*{Método de 4 pasos}
	
	Si lo que medimos es mayor a la longitud de onda, la fase se envuelve. Si mezclamos varias $\lambda$ podemos aumentar la resolución, pero aumenta el ruido.\\
	
	Para controlar el valor de $\alpha$ el interferómetro se ajusta ligeramente entre cada medición. Típicamente se utiliza un pizoelectrico (al modular, controlamos el cambio de la frecuencia),como el titanato de circonato de plomo (PZT). El phase stepping y el otro son modulaciones temporales, las de fourier son en la frecuencia.
	
	\section*{Algoritmo de Carré}
	Es ampliamente usado en interferometría óptica porque el paso es arbitrario. No es sensitivo al error de cambio de paso. Este algorimo es un algoritmo de 4 pasos con una cambio de fase desconocido.
	
	Las 4 intensidades se pueden escribir como:
	
	$$ I_1(x, y) =  $$
	
	Así queda como:
	
	$$\Delta \phi = \arctan\left(\frac{\sqrt{I_1 + I_2 - I_3 -I_4} \cdot \sqrt{3I_2-3I_3-I_1+I_4}}{I_2+I_3-I_1-I_4}\right)$$
	
	\section*{Algoritmo de 5 pasos Schwideer-Hariharan}
	Este método comprueba que la primer imagen es igual que el quinto, ya que los desplazamientos de fase es de $\pi/2 = 90^{\circ}$.
	
	Así haciendo el álgebra llegamos a:
	
	
	\subsection*{Fuentes de error}
	\begin{itemize}
		\item Desplazamiento de fase incorrecto entre las imagenes.
		\item Desfase incorrecto suele deberse a las vibraciones.
		\item Refrexiones parásitas.
		\item Error de cuantificación.
		\item No linealidad del detector.
		\item Inestabilidad de la fuente. 
	\end{itemize}
	
	
	Una vez determinada la fase a través del campo de interferencia, se puede determinar el desplazamiento correspondiente (distribución de alturas $d(x, y)$) en la superficie de prueba.
	$$ d(x, y) = \frac{\lambda}{4\pi}\phi(x, y) $$	

	
	
\end{document}